\begin{lemma}[Zero counting for a trigonometric polynomial over one period]
Let $m\in\mathbb{N}$ and let $T:\mathbb{R}\to\mathbb{R}$ be a real
trigonometric polynomial of degree at most $m$ which does not vanish
identically, i.e. there is $s\in\mathbb{R}$ with $T(s)\ne0$. Let
$t_0\in\mathbb{R}$, let $Z$ be a finite set of real numbers such that
\[
  Z\subseteq[t_0,t_0+2\pi)
  \qquad\text{and}\qquad
  T(z)=0\ \text{ for every }z\in Z,
\]
and let $t_1\in[t_0,t_0+2\pi)$ satisfy $t_1\notin Z$, $T(t_1)=0$ and
$T'(t_1)=0$, where $T'$ is the pointwise derivative of $T$. Then
\[
  \#Z+2\ \le\ 2m .
\]
\end{lemma}
\begin{proof}
Throughout, $i\in\mathbb{C}$ is the imaginary unit, and for $s\in\mathbb{R}$ we
write $\varphi(s)=e^{is}\in\mathbb{C}$.

\emph{Step 0: $T$ is differentiable.} By \noderef{TrigPolySmooth}, $T$ is
real-analytic on $\mathbb{R}$, in particular differentiable at every point; so
$T'$ is the genuine derivative of $T$ everywhere, and the hypothesis
$T'(t_1)=0$ says that the derivative of $T$ at $t_1$ vanishes.

\emph{Step 1: the associated complex polynomial.} By
\noderef{TrigPolyToComplexPoly} there is $Q\in\mathbb{C}[X]$ with
\begin{equation}\label{eq:tpzc-Q}
  \deg Q\le 2m
  \qquad\text{and}\qquad
  Q\bigl(e^{it}\bigr)=e^{imt}\,T(t)\ \text{ for every }t\in\mathbb{R}.
\end{equation}

\emph{Step 2: $Q\ne0$.} Pick $s\in\mathbb{R}$ with $T(s)\ne0$. By
\eqref{eq:tpzc-Q}, $Q(\varphi(s))=e^{ims}T(s)$. Now
$\bigl\lvert e^{ims}\bigr\rvert=1$, so $e^{ims}\ne0$, and $T(s)\ne0$ as a
complex number because the inclusion $\mathbb{R}\subseteq\mathbb{C}$ is
injective; hence $Q(\varphi(s))\ne0$ and therefore $Q$ is not the zero
polynomial. Consequently $\deg Q$ is a natural number, and \eqref{eq:tpzc-Q}
reads
\begin{equation}\label{eq:tpzc-deg}
  \deg Q\le 2m .
\end{equation}

\emph{Step 3: $\varphi$ is injective on $[t_0,t_0+2\pi)$.} Let
$s,s'\in[t_0,t_0+2\pi)$ with $e^{is}=e^{is'}$. Since $e^{w}=e^{w'}$ holds for
complex $w,w'$ exactly when $w=w'+n\,(2\pi i)$ for some $n\in\mathbb{Z}$, there
is $n\in\mathbb{Z}$ with $is=is'+2\pi i n$; dividing by $i\ne0$ gives
$s=s'+2\pi n$. From $t_0\le s<t_0+2\pi$ and $t_0\le s'<t_0+2\pi$ we get
$\lvert s-s'\rvert<2\pi$, i.e. $2\pi\lvert n\rvert<2\pi$, so $\lvert n\rvert<1$
and therefore $n=0$ and $s=s'$.

\emph{Step 4: $\#Z+1$ distinct roots of $Q$.} Put $W=Z\cup\{t_1\}$. Since
$t_1\notin Z$, $W$ is a finite set with $\#W=\#Z+1$, and
$W\subseteq[t_0,t_0+2\pi)$ by hypothesis. For $z\in Z$ we have $T(z)=0$, so
\eqref{eq:tpzc-Q} gives $Q(\varphi(z))=e^{imz}\cdot0=0$; likewise
$Q(\varphi(t_1))=e^{imt_1}\cdot0=0$ since $T(t_1)=0$. Hence every point of
$\varphi(W)$ is a root of $Q$, and by Step 3 the map $\varphi$ is injective on
$W$, so $\varphi(W)$ consists of exactly $\#Z+1$ pairwise distinct roots of $Q$.

\emph{Step 5: $w_1:=\varphi(t_1)$ is a root of $Q$ of multiplicity at least
$2$.} Consider the two $\mathbb{C}$-valued functions of a real variable
\[
  F(t)=Q\bigl(e^{it}\bigr),\qquad G(t)=e^{imt}\,T(t)\qquad(t\in\mathbb{R}).
\]
By \eqref{eq:tpzc-Q} they agree at every point, i.e. $F=G$ as functions
$\mathbb{R}\to\mathbb{C}$; consequently they have the same derivative at every
point of $\mathbb{R}$.

We compute both derivatives at an arbitrary $t\in\mathbb{R}$.
\begin{itemize}
\item The curve $\gamma(t)=e^{it}$ is differentiable as a map
$\mathbb{R}\to\mathbb{C}$ with $\gamma'(t)=i\,e^{it}$, and the polynomial
function $w\mapsto Q(w)$ is complex differentiable at every $w\in\mathbb{C}$
with complex derivative $Q'(w)$, where $Q'$ denotes the formal derivative of
$Q$ (for a polynomial the complex derivative of its associated function is the
function associated with its formal derivative). If $g:\mathbb{C}\to\mathbb{C}$
is complex differentiable at $\gamma(t)$ and $\gamma:\mathbb{R}\to\mathbb{C}$ is
differentiable at $t$, then $g\circ\gamma$ is differentiable at $t$ with
$(g\circ\gamma)'(t)=g'(\gamma(t))\,\gamma'(t)$: indeed, complex
differentiability of $g$ at $\gamma(t)$ says that the $\mathbb{R}$-linear
Fr\'echet derivative of $g$ at $\gamma(t)$ is multiplication by $g'(\gamma(t))$,
and the ordinary chain rule for Fr\'echet derivatives then gives the displayed
formula. Applying this with $g(w)=Q(w)$ yields that $F=g\circ\gamma$ is
differentiable and
\[
  F'(t)=Q'\bigl(e^{it}\bigr)\cdot i\,e^{it}.
\]
\item The curve $t\mapsto e^{imt}$ is differentiable with derivative
$t\mapsto i\,m\,e^{imt}$, and $T$ is differentiable by Step 0. By the product
rule, $G$ is differentiable and
\[
  G'(t)=i\,m\,e^{imt}\,T(t)+e^{imt}\,T'(t).
\]
\end{itemize}
Evaluating at $t=t_1$ and using $T(t_1)=0$ and $T'(t_1)=0$ gives $G'(t_1)=0$,
hence
\[
  Q'(w_1)\cdot i\,e^{it_1}=F'(t_1)=G'(t_1)=0 .
\]
Since $i\ne0$ and $\lvert e^{it_1}\rvert=1$ so $e^{it_1}\ne0$, we may cancel and
obtain
\begin{equation}\label{eq:tpzc-double}
  Q(w_1)=0\qquad\text{and}\qquad Q'(w_1)=0 .
\end{equation}
Now $Q(w_1)=0$ means $w_1$ is a root of $Q$, so by the factor theorem
$Q=(X-w_1)R$ for some $R\in\mathbb{C}[X]$. Differentiating formally,
$Q'=R+(X-w_1)R'$, so evaluating at $w_1$ gives $Q'(w_1)=R(w_1)$. By
\eqref{eq:tpzc-double} we get $R(w_1)=0$, so the factor theorem applies again
and $R=(X-w_1)S$ for some $S\in\mathbb{C}[X]$. Therefore
$Q=(X-w_1)^2S$, i.e. $(X-w_1)^2$ divides $Q$, which (as $Q\ne0$ by Step 2) says
precisely that the multiplicity of $w_1$ as a root of $Q$ is at least $2$.

\emph{Step 6: counting with multiplicity.} For $w\in\mathbb{C}$ let $\mu(w)$ be
the multiplicity of $w$ as a root of $Q$. Since $Q\ne0$ and $\mathbb{C}$ is a
field, the multiset of roots of $Q$ counted with multiplicity contains each $w$
exactly $\mu(w)$ times and has total cardinality at most $\deg Q$. Because the
elements of $\varphi(W)$ are pairwise distinct (Step 4), their contributions to
that multiset are disjoint, so
\[
  \sum_{w\in\varphi(W)}\mu(w)\ \le\ \deg Q\ \le\ 2m,
\]
the last inequality being \eqref{eq:tpzc-deg}. On the other hand, by Step 4 each
of the $\#Z$ points $\varphi(z)$, $z\in Z$, is a root of $Q$ and so has
$\mu(\varphi(z))\ge1$, while $\mu(w_1)\ge2$ by Step 5; since
$\varphi(W)=\{\varphi(z):z\in Z\}\cup\{w_1\}$ is a disjoint union of $\#Z$
points and the point $w_1$ (Step 3 and $t_1\notin Z$),
\[
  \sum_{w\in\varphi(W)}\mu(w)\ \ge\ \#Z\cdot 1+2=\#Z+2 .
\]
Combining the two displays gives $\#Z+2\le 2m$, as claimed.
\end{proof}
